%% ============================================================
%%  Frontend Reference — Paper Trading Dashboard
%%  Compile: pdflatex frontend_report.tex  (twice for refs)
%% ============================================================
\documentclass[12pt, a4paper]{article}

\usepackage{amsmath}
\usepackage{booktabs}
\usepackage{geometry}
\usepackage{hyperref}
\usepackage{xcolor}
\usepackage{array}
\usepackage{enumitem}
\usepackage{caption}
\usepackage{float}
\usepackage{parskip}
\usepackage{listings}
\usepackage{fancyvrb}

\geometry{margin=2.5cm}

\hypersetup{colorlinks=true, linkcolor=blue!60!black, urlcolor=blue!60!black}

\definecolor{codebg}{HTML}{0d0d14}
\definecolor{codeframe}{HTML}{1e1e2e}

\lstset{
  basicstyle=\ttfamily\small,
  backgroundcolor=\color{codebg!10},
  frame=single,
  framerule=0.4pt,
  rulecolor=\color{codeframe},
  breaklines=true,
  captionpos=b,
  showstringspaces=false,
}

\title{\textbf{Frontend Reference}\\[0.5em]
\large Paper Trading Dashboard — Architecture and Page Guide}
\author{}
\date{May 2026}

\begin{document}
\maketitle
\tableofcontents
\newpage


%% ============================================================
\section{Overview}
%% ============================================================

The paper trading dashboard is a React/TypeScript single-page application (SPA)
that provides a real-time window into the N3 DVOL strategy system running in the
FastAPI backend. Every page fetches data directly from the backend REST API; there
is no intermediate state-management library (no Redux, no Zustand) — each page
owns its own local state with React's built-in \texttt{useState} and \texttt{useEffect}.

The dashboard was built for a single operator running a small number of strategies
on BTCUSDT perpetual futures. It is not designed for multiple simultaneous users.
All destructive actions (kill switch, strategy demotion) are protected by the
backend's own validation; the frontend is a thin presentation layer.


%% ============================================================
\section{Technology Stack}
%% ============================================================

\begin{table}[H]
\centering
\begin{tabular}{ll}
\toprule
Concern & Library / Version \\
\midrule
Framework       & React 19 with TypeScript 4.9 \\
Routing         & React Router DOM v7 \\
Charts          & Recharts 3 \\
Icons           & Lucide React \\
Build tool      & Create React App (\texttt{react-scripts 5}) \\
Dev port        & 3000 (local), 80 (Docker via nginx) \\
API base URL    & \texttt{http://localhost:8000/api} (configurable via \texttt{REACT\_APP\_API\_URL}) \\
\bottomrule
\end{tabular}
\end{table}

There are no UI component libraries (no Material UI, no Chakra). All styling is
hand-written CSS in \texttt{App.css} using a consistent dark-theme design system
with class names like \texttt{.card}, \texttt{.badge}, \texttt{.table-wrapper},
\texttt{.positive}, and \texttt{.negative}.


%% ============================================================
\section{File Structure}
%% ============================================================

\begin{Verbatim}[fontsize=\small]
paper_trading/frontend/src/
│
├── App.tsx             Navigation bar + route definitions
├── App.css             Global dark-theme design system
├── index.tsx           React entry point
│
├── api/
│   ├── client.ts       apiFetch / apiPost / apiPostBody / apiPatch / apiDelete
│   ├── index.ts        Named api.* methods for every endpoint
│   └── types.ts        TypeScript interfaces for all API responses
│
├── components/
│   └── BotConfigPanel.tsx   Slide-out gear panel for market/strategy config
│
└── pages/
    ├── Dashboard.tsx        Live market data, signal state, position, equity
    ├── Signals.tsx          Daily signal evaluation log
    ├── Trades.tsx           Full trade journal
    ├── TradeDetail.tsx      Per-trade PnL breakdown with attribution
    ├── Performance.tsx      Equity curve, Sharpe, drawdown chart, yearly table
    ├── Replay.tsx           Historical backtest certification (not live trades)
    ├── ForwardLog.tsx       P3 shadow strategy independence monitor
    ├── RiskDashboard.tsx    Portfolio exposure, limit gauges, drawdowns
    ├── StrategyPipeline.tsx Strategy lifecycle manager (research → validated)
    ├── Alerts.tsx           Inbox for system events with category filter
    ├── DataQuality.tsx      Feed freshness and completeness checks
    ├── Experiments.tsx      Research run log with metrics drilldown
    └── SystemHealth.tsx     Scheduler status, data freshness, error logs
\end{Verbatim}


%% ============================================================
\section{API Layer}
%% ============================================================

\subsection{client.ts — HTTP primitives}

All HTTP calls flow through four functions in \texttt{client.ts}:

\begin{Verbatim}[fontsize=\small]
const BASE = process.env.REACT_APP_API_URL ?? "http://localhost:8000/api";

apiFetch<T>(path, opts?)      // GET (default) or any method
apiPost<T>(path)              // POST with no body
apiPostBody<T>(path, body)    // POST with JSON body
apiPatch<T>(path, body)       // PATCH with JSON body
apiDelete<T>(path)            // DELETE
\end{Verbatim}

Every function throws if the response is not \texttt{2xx}. Pages catch the
exception and store the message in local \texttt{error} state, displaying it
inline with red text rather than crashing.

\subsection{index.ts — Named API methods}

\texttt{api.*} provides one named method per endpoint, typed against the
interfaces in \texttt{types.ts}. All pages import \texttt{api} from this
module and never call \texttt{apiFetch} directly. This ensures a single
place to update if an endpoint path changes.

\begin{table}[H]
\centering\small
\begin{tabular}{ll}
\toprule
Method & Endpoint \\
\midrule
\texttt{api.dashboard()}            & \texttt{GET /api/dashboard} \\
\texttt{api.signalHistory(n)}       & \texttt{GET /api/signals/history?limit=n} \\
\texttt{api.trades()}               & \texttt{GET /api/trades} \\
\texttt{api.trade(id)}              & \texttt{GET /api/trades/:id} \\
\texttt{api.performance()}          & \texttt{GET /api/performance} \\
\texttt{api.replay(start, train)}   & \texttt{GET /api/replay?start=\&include\_train=} \\
\texttt{api.forwardLog(strategy)}   & \texttt{GET /api/forward-log/p3?strategy=} \\
\texttt{api.riskState()}            & \texttt{GET /api/risk/state} \\
\texttt{api.strategies()}           & \texttt{GET /api/strategies} \\
\texttt{api.strategyUpdateStatus()} & \texttt{POST /api/strategies/:name/status} \\
\texttt{api.alerts(n)}              & \texttt{GET /api/alerts?limit=n} \\
\texttt{api.alertRead(id)}          & \texttt{POST /api/alerts/:id/read} \\
\texttt{api.dataQuality()}          & \texttt{GET /api/system/data-quality} \\
\texttt{api.experiments()}          & \texttt{GET /api/experiments} \\
\texttt{api.systemHealth()}         & \texttt{GET /api/system/health} \\
\texttt{api.runDailySignal()}       & \texttt{POST /api/jobs/run-daily-signal} \\
\texttt{api.checkExits()}           & \texttt{POST /api/jobs/check-exits} \\
\texttt{api.configAvailable()}      & \texttt{GET /api/config/available} \\
\texttt{api.configActive()}         & \texttt{GET /api/config/active} \\
\texttt{api.configAdd()}            & \texttt{POST /api/config/active} \\
\texttt{api.configRemove()}         & \texttt{DELETE /api/config/active/:market/:strategy} \\
\bottomrule
\end{tabular}
\end{table}

\subsection{types.ts — TypeScript interfaces}

Every API response has a corresponding TypeScript interface. The key ones are:

\begin{itemize}[noitemsep]
  \item \textbf{DashboardData} — combined market + signal + position + equity snapshot
  \item \textbf{SignalRecord} — one row per daily evaluation (dvol, n3\_z, entry\_signal, reason)
  \item \textbf{TradeRecord} — full trade journal row with PnL attribution fields
  \item \textbf{PerformanceData} — aggregate stats, equity history, yearly breakdown
  \item \textbf{PortfolioState} — live exposure, limit usage, open positions, drawdowns
  \item \textbf{StrategyPipelineRow} — lifecycle status + live stats per strategy
  \item \textbf{ForwardLogResponse} — P3 shadow evaluation log + independence monitor stats
  \item \textbf{AlertRecord} — event notification with category, title, body, read state
  \item \textbf{DataQualityReport} — list of named feed checks with ok/warn/error status
  \item \textbf{ExperimentRun} — research run metadata: script, params, metrics, verdict
\end{itemize}


%% ============================================================
\section{Navigation and Routing}
%% ============================================================

\texttt{App.tsx} defines the navigation bar and the route table. The nav bar
is a horizontal strip that stays fixed at the top. Each link is a
\texttt{NavLink} from React Router and gets an \texttt{active} CSS class when
its route matches. At the right end of the nav bar sits the \texttt{BotConfigPanel}
gear icon (Section~\ref{sec:botconfig}).

\begin{table}[H]
\centering\small
\begin{tabular}{lll}
\toprule
Route & Page & Primary purpose \\
\midrule
\texttt{/}              & Dashboard         & Live system state \\
\texttt{/signals}       & Signals           & Daily evaluation log \\
\texttt{/trades}        & Trades            & Trade journal \\
\texttt{/trades/:id}    & TradeDetail       & Per-trade PnL attribution \\
\texttt{/performance}   & Performance       & Equity curve \& aggregate stats \\
\texttt{/replay}        & Replay            & Research backtest certification \\
\texttt{/forward-log}   & ForwardLog        & P3 shadow strategy log \\
\texttt{/risk}          & RiskDashboard     & Portfolio exposure \& limits \\
\texttt{/pipeline}      & StrategyPipeline  & Strategy lifecycle manager \\
\texttt{/alerts}        & Alerts            & System event inbox \\
\texttt{/data-quality}  & DataQuality       & Feed health checks \\
\texttt{/experiments}   & Experiments       & Research run log \\
\texttt{/health}        & SystemHealth      & Scheduler \& data freshness \\
\bottomrule
\end{tabular}
\end{table}


%% ============================================================
\section{Page Reference}
%% ============================================================

\subsection{Dashboard (\texttt{/})}
\label{sec:dashboard}

The landing page. It is the only page with an automatic refresh loop — it
calls \texttt{GET /api/dashboard} every 30 seconds via \texttt{setInterval}
and clears the interval on unmount.

\paragraph{Sections.}
\begin{enumerate}[noitemsep]
  \item \textbf{Market} — BTC spot price, current 8h funding rate, last update time.
  \item \textbf{Signal} — DVOL level (coloured green if $\geq 54$), N3 z-score,
        DVOL filter pass/fail badge, signal outcome (LONG or FLAT),
        and the human-readable reason string from the last evaluation.
  \item \textbf{Position} — if a trade is open: entry price, unrealised PnL in basis
        points, and a countdown showing hours until the 24h exit. If flat: a grey
        ``no open position'' card.
  \item \textbf{Equity} — running equity (starts at \$10,000), realised PnL as a
        percentage, and current drawdown from peak.
\end{enumerate}

\paragraph{Manual trigger.}
The ``Run Signal Now'' button calls \texttt{POST /api/jobs/run-daily-signal},
then immediately refreshes the dashboard. This is the mechanism for testing the
signal evaluator without waiting for midnight UTC.

\paragraph{Design notes.}
DVOL is coloured \textcolor[HTML]{4ade80}{green} when the regime filter passes
($\geq 54$) and \textcolor[HTML]{f87171}{red} when it fails. The LONG/FLAT
badge uses the same green/grey convention. This makes the signal state readable
at a glance without reading numbers.

---

\subsection{Signals (\texttt{/signals})}

Shows every daily N3 evaluation stored in the database, newest first (the
list from the API is oldest-first; the page reverses it for display).

\paragraph{Summary cards.}
Three counters at the top: total evaluations, days where DVOL $\geq 54$
(the regime filter), and days where the full entry signal fired. The
``DVOL Filter Pass'' counter as a fraction of total evaluations shows
what proportion of days the market was in the high-vol regime.

\paragraph{Table columns.}
Date, DVOL, 30d mean, 30d std, N3 z-score (highlighted green if $> 0.75$),
DVOL filter pass/fail badge, signal badge (LONG or FLAT), and the full
reason string from the evaluator.

\paragraph{Empty state.}
If no signals exist yet (fresh system), the table is replaced with a prompt
to run the daily signal job.

---

\subsection{Trades (\texttt{/trades})}

Fetches all trades from \texttt{GET /api/trades} and renders them in a
scrollable table. Each row links to the corresponding \texttt{/trades/:id}
detail page.

Columns: trade ID, strategy, market, side (LONG/SHORT badge), entry
timestamp, entry price, planned exit, status (OPEN/CLOSED badge), and
net PnL in basis points coloured green/red. Open trades show ``---'' for
the PnL column since the position has not yet been closed.

---

\subsection{TradeDetail (\texttt{/trades/:id})}

Per-trade PnL attribution. Fetches \texttt{GET /api/trades/:id} for the
single record. Displays:

\begin{itemize}[noitemsep]
  \item Entry and exit prices, timestamps, and the human-readable entry reason.
  \item PnL decomposition: gross price return (log-return, in bp), funding PnL
        (three settlements over 24h, negative when longs pay), fees (6bp
        round-trip, maker cost model), and net PnL as the sum.
  \item N3 z-score and DVOL level at entry, for context on why the signal fired.
\end{itemize}

---

\subsection{Performance (\texttt{/performance})}

Shows aggregate performance of all closed trades. Requests
\texttt{GET /api/performance} which returns pre-computed statistics
plus the full equity curve history.

\paragraph{Summary cards.}
Sharpe (annualised; highlighted green if $> 2$), total PnL in bp, max drawdown,
win rate, average winning trade (bp), average losing trade (bp), and profit factor.

\paragraph{Charts.}
Two Recharts \texttt{LineChart} components below the cards:
\begin{enumerate}[noitemsep]
  \item \textbf{Equity curve} — running equity in USD against time, with a dashed
        reference line at \$10,000 (the starting capital). Points are individual
        equity curve rows from the database, not interpolated.
  \item \textbf{Drawdown} — percentage drawdown from peak equity against the same
        timeline, in red. Both charts use the same \texttt{ResponsiveContainer}
        wrapper so they resize with the browser window.
\end{enumerate}

\paragraph{Year-by-year table.}
If the backend returns a \texttt{yearly\_breakdown} array, a third section
shows trades, net PnL, Sharpe, and win rate per calendar year. This is
useful for spotting regime-driven variance.

\paragraph{Empty state.}
When there are no closed trades yet, the entire page is replaced with a
single card saying ``No closed trades yet.'' The charts and table are not
rendered.

---

\subsection{Replay (\texttt{/replay})}

\textbf{Important:} this page does not show live paper trading results.
It runs the historical research backtest from scratch against the raw parquet
data files (\texttt{data/raw/}) and verifies that the live backend code
produces numbers within 15\% of the research report targets. It is a
certification tool.

\paragraph{How it works.}
Clicking ``Run Replay'' calls \texttt{GET /api/replay?start=2024-01-01}.
The backend reads the 1-minute kline and hourly DVOL parquet files, computes
the 30d rolling z-score, samples daily, applies the frozen rule
(\texttt{N3z > 0.75 AND DVOL} $\geq 54$), and returns per-trade results with
cumulative PnL. The process takes 5--10 seconds because it scans 1.7M bars.

\paragraph{Match badges.}
Each summary card (Trades, Sharpe, Net PnL, Max Drawdown, Win Rate) shows a
green \textbf{MATCH} or red \textbf{DIFF} badge computed as:
\[
  \text{pass} \iff \frac{|\text{actual} - \text{target}|}{|\text{target}|} \leq 0.15
\]
The 15\% tolerance accounts for minor differences in data alignment
between the research scripts and the live data pipeline.

\paragraph{Reference targets.}
\begin{center}
\begin{tabular}{lr}
\toprule
Metric & Research target \\
\midrule
Trades       & 199 \\
Sharpe       & $+2.95$ \\
Net PnL      & $+11{,}228$ bp \\
Max drawdown & $-1{,}612$ bp \\
Win rate     & $53.8\%$ \\
\bottomrule
\end{tabular}
\end{center}

\paragraph{Checkbox: include 2023.}
When ticked, the replay also shows 2023 training-period results. These are
displayed at reduced opacity in the period breakdown table to distinguish
in-sample from out-of-sample rows.

\paragraph{Error state.}
If the parquet files are not present (excluded from the repository),
the backend returns HTTP 503 and the frontend shows a specific message:
``Parquet data files not found. Run data/download.py\ldots''

---

\subsection{P3 Log / Forward Log (\texttt{/forward-log})}

The shadow monitoring page for the P3 OI-Price Divergence strategy, which
is in \texttt{shadow} status and under forward review until $\geq$\,2026-08-14.

\paragraph{Purpose.}
Every day P3 is evaluated alongside N3, regardless of whether a trade is
opened. This page records every evaluation so that the independence of P3
from N3 can be measured in real time. The key metric is the Sharpe on
\emph{exclusive} P3 trades --- days where P3 fired but N3 did not.
OOS research baseline: Sharpe $= +5.18$ on 57 exclusive trades ($p = 0.007$).

\paragraph{Summary stat boxes.}
Three side-by-side boxes: ``All trades'', ``P3-exclusive (independence
monitor)'', and ``N3 overlap''. Each shows $n$, Sharpe, total PnL, and win
rate. The exclusive-trade box is the forward review's primary metric.

\paragraph{Evaluation counter.}
A fourth card shows total evaluations, traded days, exclusive days (P3
fired, N3 did not), and overlap days (both strategies fired).

\paragraph{Row toggle.}
By default only rows where the signal fired are shown. The ``Show all
evaluations'' button reveals no-signal rows at 50\% opacity. No-signal rows
are important because they prove the rule stayed frozen: if P3 never fires
during a long low-volatility stretch, the log still records an evaluation
every day confirming the evaluator is running correctly.

\paragraph{Table columns.}
Date, DVOL, regime (DD/DU/NN), $\Delta$P 24h (price change, red if negative),
$\Delta$OI 24h (open interest change, red if negative), signal (YES/no badge),
exclusive/overlap badge, N3 also fired (blue checkmark), trade status, and
net PnL in basis points.

\paragraph{Empty state.}
When no P3 signals exist yet (fixed by the bug described in the previous
section), the page shows ``No rows to display'' rather than crashing.

---

\subsection{Portfolio Risk (\texttt{/risk})}

Live exposure dashboard. Fetches \texttt{GET /api/risk/state}, which returns
current open positions, daily PnL, and limit configuration.

\paragraph{Limit gauges.}
Two horizontal bar charts at the top:
\begin{itemize}[noitemsep]
  \item \textbf{Open positions} — $n_{\text{open}} / 3$ (max 3 total positions).
        Bar turns yellow at $50\%$ usage, red at $80\%$.
  \item \textbf{Daily loss} — current daily PnL against the $-500$ bp limit.
        Only fills when the day is in loss.
\end{itemize}

\paragraph{Hard limits table.}
Shows the current value, configured maximum, and an OK/BREACH status badge
for total open positions and daily PnL.

\paragraph{Open positions table.}
Lists every currently open trade: trade ID, strategy, market, side badge,
entry price, DVOL at entry, and planned exit time.

\paragraph{Strategy trailing drawdowns.}
For each strategy with at least 5 trades, shows the trailing drawdown on
the last 20 trades against the $-20\%$ limit. Strategies that breach this
are shown as ``PAUSED'' (the backend will block new entries for them).

---

\subsection{Strategy Pipeline (\texttt{/pipeline})}

Displays every known strategy and its lifecycle status. The formal stages are:

\begin{center}
research $\to$ candidate $\to$ shadow $\to$ validated $\to$ paused / killed
\end{center}

Each status has a distinct colour: grey (research), blue (candidate), yellow
(shadow), green (validated), orange (paused), red (killed).

\paragraph{Status counts.}
A row of small cards at the top shows how many strategies are in each stage.
Only stages with at least one strategy appear.

\paragraph{Per-strategy row.}
Strategy name (monospace), status badge, description string from the backend,
number of live evaluations, number of trades opened, last promotion date, and
a note field. Each row has an ``Edit'' button that expands an inline form for
changing the status and adding a note.

\paragraph{Editing.}
The inline edit form offers a dropdown of all valid statuses and a free-text
note field. Clicking ``Save'' calls \texttt{POST /api/strategies/:name/status}
and refreshes the table. This is the primary mechanism for promoting P3 from
shadow to validated (or killing it) after the forward review.

---

\subsection{Alerts (\texttt{/alerts})}

An inbox for system events. Fetches both \texttt{GET /api/alerts} (last 200
records) and \texttt{GET /api/alerts/summary} (unread counts by category).

\paragraph{Alert categories.}
\begin{table}[H]
\centering\small
\begin{tabular}{ll}
\toprule
Category & When it fires \\
\midrule
\texttt{signal\_fired}     & An entry signal was evaluated as LONG and a trade was opened \\
\texttt{trade\_closed}     & A trade reached its 24h expiry and was closed \\
\texttt{data\_failed}      & Binance or Deribit API returned an error \\
\texttt{scheduler\_missed} & The daily job did not fire at midnight UTC \\
\texttt{risk\_blocked}     & A portfolio risk limit prevented a trade from opening \\
\bottomrule
\end{tabular}
\end{table}

\paragraph{Unread indicators.}
Unread alerts have a small blue dot in the leftmost column and full opacity.
Read alerts are at 55\% opacity. A ``Mark all read'' button appears when
there are unread alerts.

\paragraph{Category filter.}
A row of filter buttons (all / signal\_fired / trade\_closed / \ldots) lets
the operator focus on a single alert type. Clicking a category summary card
at the top also filters to that category.

\paragraph{Expandable detail.}
Clicking any row expands a full-width panel below it showing the alert
body text (pre-formatted) and the market name. This is where the
per-trade PnL summary, the data failure error message, or the risk block
reason is displayed.

---

\subsection{Data Quality (\texttt{/data-quality})}

Fetches \texttt{GET /api/system/data-quality} which runs a set of named
checks against the latest stored market data.

\paragraph{Summary cards.}
Overall status (OK / WARN / ERROR), total check count, warning count, and
error count. The overall status is the worst of any individual check.

\paragraph{Checks table.}
Each check row shows: check name (monospace), human-readable detail, current
value, expected value or threshold, and a coloured status badge. Examples of
checks include: ``last DVOL update within 2 hours'', ``BTC price within 5\%
of mark price'', ``funding rate within expected range''.

\paragraph{Incident use.}
This page is the first place to look during a data staleness alert
(Severity S2). A check with status \texttt{error} directly identifies
which feed is stale or unreachable.

---

\subsection{Experiments (\texttt{/experiments})}

A read-only log of every research run recorded in the \texttt{experiment\_runs}
database table. Fetches \texttt{GET /api/experiments}.

\paragraph{Verdict summary.}
A row of small cards counts runs by verdict: passed, failed, killed, pending.

\paragraph{Runs table.}
Columns: run ID (unique slug like \texttt{25\_du\_short\_20260514}), strategy
name, script name, data range (start → end), short commit hash (8 chars), verdict
badge, and creation date.

\paragraph{Expandable drilldown.}
Clicking a row expands an inline panel showing three sections:
\begin{enumerate}[noitemsep]
  \item \textbf{Metrics} — key-value pills (e.g.\ \texttt{sharpe: -0.44},
        \texttt{n\_trades: 183}).
  \item \textbf{Parameters} — the configuration of the research run
        (thresholds, data ranges, cost model).
  \item \textbf{Notes} — free-text verdict rationale.
\end{enumerate}

This page gives an auditable record of every hypothesis tested and its
outcome, supporting the ``one kill attempt per signal'' discipline.

---

\subsection{System Health (\texttt{/health})}

Fetches \texttt{GET /api/system/health} for the system status snapshot and
\texttt{GET /api/system/logs} for recent log lines.

Displays: overall health status (healthy/degraded), last Binance data update,
market data staleness flag, last signal computation time, signal staleness flag,
open position count, stale threshold in minutes, and database status.

The log table shows recent \texttt{system\_logs} rows: timestamp, level
(INFO/WARNING/ERROR), component name, and message. Error-level rows are
highlighted red.

---

%% ============================================================
\section{BotConfigPanel}
\label{sec:botconfig}
%% ============================================================

The gear icon in the top-right of the navigation bar opens the
\texttt{BotConfigPanel}, a slide-in panel from the right edge of the screen.
It is a component (not a page) embedded directly in the nav bar, so it is
accessible from any page without navigation.

\paragraph{Active bots list.}
The top section shows every active \texttt{BotConfig} row: the market
icon, strategy display name, market identifier, and a remove button (×).
Removing a bot calls \texttt{DELETE /api/config/active/:market/:strategy}
and refreshes the list. An active bot is the driver for the daily signal
job --- removing it stops future signal evaluations for that
market/strategy pair.

\paragraph{Add new bot.}
A two-step flow below the active list:
\begin{enumerate}[noitemsep]
  \item \textbf{Select market} — buttons rendered from
        \texttt{GET /api/config/available}. Currently only BTCUSDT is listed.
  \item \textbf{Select strategy} — after a market is chosen, only
        strategies compatible with that market are shown. Each strategy card
        displays its status badge, description, tags, and hold period.
        Already-active strategies are shown at 50\% opacity and cannot be
        selected again. Strategies with status \texttt{coming\_soon} are
        similarly disabled.
\end{enumerate}
Clicking ``+ Add Bot'' calls \texttt{POST /api/config/active} and refreshes
both lists.

\paragraph{Backdrop.}
A dark translucent overlay covers the rest of the page when the panel is open.
Clicking it closes the panel without saving.


%% ============================================================
\section{State Management Pattern}
%% ============================================================

The frontend uses no global state library. Every page manages its own state
with \texttt{useState} and fetches data in \texttt{useEffect}. The pattern
used throughout is:

\begin{Verbatim}[fontsize=\small]
const [data,    setData]    = useState<T | null>(null);
const [loading, setLoading] = useState(true);
const [error,   setError]   = useState<string | null>(null);

const load = () => {
  setLoading(true);
  api.someEndpoint()
    .then(setData)
    .catch(e => setError(e.message))
    .finally(() => setLoading(false));
};

useEffect(() => { load(); }, []);
\end{Verbatim}

Pages that support manual refresh expose a \texttt{load} function and attach it
to a ``$\circlearrowleft$ Refresh'' button. The Dashboard is the only page that
also sets up a \texttt{setInterval(load, 30000)} for automatic 30-second polling.

Pages that support write operations (StrategyPipeline, Alerts, BotConfigPanel)
call the write method, then call \texttt{load()} / \texttt{refresh()} to refetch
and display the updated state.


%% ============================================================
\section{Design System}
%% ============================================================

All visual styling is in \texttt{App.css}. The design follows a dark theme
with a near-black background (\texttt{\#0d0d17}) and a slightly lighter surface
colour (\texttt{\#111118}) for cards.

\paragraph{CSS class conventions.}
\begin{itemize}[noitemsep]
  \item \textbf{\texttt{.card}} — a bordered rounded box with 16px padding.
        Used for every metric tile and content section.
  \item \textbf{\texttt{.cards-grid}} — a CSS grid that auto-fills cards
        at a minimum width of 160px, wrapping on narrow viewports.
  \item \textbf{\texttt{.badge}} — a small pill label. Variants:
        \texttt{badge-green}, \texttt{badge-red}, \texttt{badge-yellow},
        \texttt{badge-blue}, \texttt{badge-gray}.
  \item \textbf{\texttt{.positive}} / \textbf{\texttt{.negative}} — green
        (\texttt{\#4ade80}) and red (\texttt{\#f87171}) text for PnL values.
  \item \textbf{\texttt{.table-wrapper}} — a horizontally scrollable container
        wrapping a \texttt{<table>}, with alternating row shading.
  \item \textbf{\texttt{.chart-container}} — card wrapper for Recharts charts
        with a grey title and consistent padding.
  \item \textbf{\texttt{.loading}} — centred grey italic text for loading states.
  \item \textbf{\texttt{.btn}} — base button style. Variants:
        \texttt{btn-primary} (indigo fill), \texttt{btn-secondary}
        (dark fill), \texttt{btn-ghost} (transparent, border only).
\end{itemize}

\paragraph{Colour semantics.}
\begin{table}[H]
\centering\small
\begin{tabular}{lll}
\toprule
Colour & Hex & Meaning \\
\midrule
Green  & \texttt{\#4ade80} & Positive PnL, filter pass, OK status \\
Red    & \texttt{\#f87171} & Negative PnL, filter fail, error status \\
Yellow & \texttt{\#fbbf24} & Warning, shadow strategy, elevated limits \\
Blue   & \texttt{\#60a5fa} & N3 co-fire indicator, trade-closed alerts \\
Orange & \texttt{\#f97316} & Paused strategy, risk-blocked alerts \\
Indigo & \texttt{\#6366f1} & Primary action buttons, equity curve line \\
Slate  & \texttt{\#64748b} & Muted text, secondary labels \\
\bottomrule
\end{tabular}
\end{table}


%% ============================================================
\section{Refresh Behaviour Summary}
%% ============================================================

\begin{table}[H]
\centering\small
\begin{tabular}{lll}
\toprule
Page & Auto-refresh & Manual refresh \\
\midrule
Dashboard         & Every 30 seconds & ↻ button + ``Run Signal Now'' \\
Signals           & No               & Page reload \\
Trades            & No               & Page reload \\
TradeDetail       & No               & Page reload \\
Performance       & No               & Page reload \\
Replay            & No (on demand)   & ``Run Replay'' button \\
P3 Log            & No               & ↻ button \\
Risk              & No               & ↻ button \\
Pipeline          & No               & ↻ button (post-edit) \\
Alerts            & No               & ↻ button + ``Mark all read'' \\
Data Quality      & No               & ↻ button \\
Experiments       & No               & ↻ button \\
System Health     & No               & Page reload \\
\bottomrule
\end{tabular}
\end{table}

Only the Dashboard auto-polls. All other pages are static once loaded.
For near-real-time monitoring after a manual signal trigger, use the Dashboard.


%% ============================================================
\section{Common Misconceptions}
%% ============================================================

\paragraph{Replay $\neq$ Performance.}
The Replay page runs a historical simulation against parquet files and shows
what would have happened. The Performance page shows what actually happened
in the live paper trading system. On a fresh system with no live trades,
Performance shows ``No closed trades yet'' while Replay shows the full
historical backtest results. These are two separate data sources.

\paragraph{P3 Log shows no-signal days.}
Every daily evaluation is recorded in the forward log, even days where the
signal did not fire. Toggling ``Show all evaluations'' reveals these rows.
This is by design: no-signal days prove the evaluator is running and the rule
has not been changed.

\paragraph{The Kill Switch is not on this dashboard.}
There is no kill switch toggle visible on the Risk page frontend in the
current implementation. The kill switch must be operated via the API directly
(see \texttt{docs/runbook.md} Section 4) or via the command line. The Risk
page does show the portfolio limits and exposure, not the kill switch state.

\paragraph{BotConfigPanel controls what runs, not what is displayed.}
Removing a bot from the config panel stops future evaluations for that
market/strategy pair. It does not delete historical signals or trades. Adding
a bot starts evaluations from the next daily job at midnight UTC.

\end{document}
